\section{Unproved extensions: exact status and current gates}
\label{sec:unproved}

This part contains questions and proposed extensions, not additional
theorems. The transfers in Parts~I--II have proofs and recorded checks.
An analytic hypothesis absent from the stipulated theorem is not thereby
proved necessary, and a failed reduction is not a logical independence
result. The following register prevents those distinctions from being
lost when the conclusions are reused.

\paragraph{CX-U01: an unrestricted Dirichlet-only HN transfer.}
The checked general route reaches RH for the Dedekind zeta function of
the bounded solution field (CX-K01). Applying prime-ideal counting to
the field itself avoids the unnecessary Galois-closure gate.
The bounded abelian-witness promise is proved in CX-K03.
The cubic in CX-K04 refutes replacement of every witness by an abelian
one; CX-R02 separately refutes a virtual rank-one representation
decomposition. Neither refutation excludes a different complexity
proof. A new argument relating rational Dirichlet zeros to the needed
prime supply, or a different AM protocol, would be required to go
further. The unconditional prime-ideal estimate inspected in
Lagarias--Odlyzko does not keep the primes at polynomial bit length
with the available solution-field bounds.

\paragraph{CX-U02: an unrestricted Dirichlet-only factoring transfer.}
The main hypothesis classification is proved: CX-N15 gives a complete
general factoring algorithm under $M$, using precisely the sufficient
family $\zeta_{\mathbb Q}$ and
$\zeta_{\mathbb Q(\zeta_{2\nu},p^{1/\nu})}$ for its explicit construction
degrees $\nu\le d^2$. It constructs the maximal coefficient Teichmüller
unit by a Frobenius eigenspace and equivariant interpolation, without
the historical radical oracle or the unretrieved internal EVD94 reduction.
Those retrieval limitations remain historical source facts, not gaps in
CX-N15. CX-N13 and CX-N14 give smaller-hypothesis cases.

The stronger extension asking for a general proof from $A_D$ alone
remains unproved. A reduction to polynomials split over the prime field
does not itself remove the root-of-unity extensions used to split
automorphisms. After CX-N14 the uncovered branch can be restricted to
$p>d$, $r\nmid p(p-1)$ and
$r^2\mid p^{\operatorname{ord}_r(p)}-1$.
A coefficient Teichmüller unit of maximal $r$-power order is a proved
sufficient extra input there. CX-N15 supplies it under its explicit
Dedekind-zeta family; no reduction of that construction to rational
Dirichlet characters, and no impossibility of such a reduction, is proved.
One further natural shortcut was checked directly: BGL Theorem~3
requires roots for every prime divisor of $\Phi_k(p)$ and retains the
largest such prime in its running time, even for one desired nonresidue.
Its alternative Theorem~(5.1) requires a degree-$r$ irreducible polynomial
when $r\nmid k$, exactly the missing input at $k=\operatorname{ord}_r(p)<r$;
see \cite[\S5, p.~20 and algorithm pp.~20--24]{CXNBGL01}.
Thus that source-based substitution does not close the smaller gate.

\paragraph{CX-U03: general polynomial-time finite-field factoring.}
The Evdokimov bound in CX-N09 and the independent composition CX-N15
are quasipolynomial in the degree.
The special-field and structural-parameter results in CX-N10 preserve
their parameters. Deleting an analytic hypothesis changes none of
these exponents. No general deterministic polynomial-time factoring
claim is made. The source-specific transfers answer the requested
hypothesis question without asserting this stronger algorithm.

\paragraph{CX-U04: unrestricted class-group search complexity.}
CX-N07 proves the factor-base transfer, and CX-N08 proves what the
index certificate establishes. The cost of finding enough independent
relations remains an algorithm-specific question. CX-N12 supplies a
concrete rigorous imaginary-quadratic computation theorem with its
randomness and subexponential bound retained. No general smoothness
heuristic is promoted to a theorem by substituting the GRH premise.

\paragraph{CX-U05: equations, symbolic circuits, and GCT separation.}
For an optimal-cardinality succinct h.s.o.p., the historical EXPH
proof still needs the equation-access assumption of GCT~V
Theorem~5.5. A polynomial-size separating system has a different,
already unconditional PSPACE construction. The ETH counter and the
Boolean circuit lower bound supply neither symbolic permanent
circuits nor graded equivariant coordinate-ring multiplicities.
The proved rank dictionary loses information, as CX-G05 demonstrates.
The remaining gates are therefore the specified equation algorithm,
arithmetic conversion, or obstruction construction. No conclusion
about P versus NP, VP versus VNP, or general determinant border
complexity is claimed beyond the explicit statements in Part~I.

\paragraph{CX-U06: stronger quantitative ETH transfers.}
The source theorem fixes clause width. The paper does not interchange
that quantifier with a width growing with input length. A graph's
vertex parameter is not its treewidth or solution-size parameter.
Likewise the Circuit-SAT encoding introduces one variable per gate,
and the nearly linear Boolean lower bound is in
$\mathrm E^{\mathrm{NP}}$. The exact reduction size, solver prefactor,
and computational model must be supplied before any stronger
fine-grained or parameterized conclusion follows.

\section{Source and verification conventions}
\label{sec:source-conventions}

The accompanying \texttt{QUESTIONS.tsv} records each treated question
and its proved finding, refutation, or remaining gate.
\texttt{CHECK\_EVENTS.tsv} records actual checks using the categories
specified in the brief. A targeted proof check means that the stated
proof or the relevant hypothesis-bearing part was read; it does not
mean a formal certificate. A cited theorem whose statement is used
without an independent reproof is identified as such.

The research run is \texttt{CX\_GRH-20261121}, begun at
23:00:34 UTC on 20 November 2026, or 00:00:34 on 21 November in
Europe/Berlin. The coordinator and three researchers used
GPT-6 Astra Pro. Their effort setting was not exposed by the runtime
and is not guessed. Dates in the consolidated registers use the
Europe/Berlin calendar; original UTC dates in component working notes
are explicitly identified. No new finding was checked by a Lean
kernel in this run.

The prior run's \texttt{BRIDGES\_QUESTIONS.tsv} rows Q07 and Q08
were read as earlier work, then rechecked: Q07's broad Dedekind gate
was sharpened to the bounded solution field, and Q08's Dirichlet
scope was verified against Bach's original Theorems~2--4 and
Schoof's complete quadratic-character proof. Their earlier check
labels are not inherited as certificates for this paper.

Archive scan numbers count the cover as scan~1. The paper also gives
the manuscript's sheet number, when present. Mathematical content
was inspected on the cited scans. These historical references locate
constructions and comparisons; none is cited as a manuscript proof
of a modern complexity statement. In particular the crossed-out
passage in cote~34 retains that qualification.

The attached sources remain the primary record for the stipulated
2026 theorems. Literature links in the bibliography identify the
consulted editions or author versions. Differences between printed
journal pagination and author-PDF pagination are stated at the
point of citation. Retrieval failures and source limitations are
recorded without inventing statement numbers or attributing a
full-proof audit to an abstract.
